\documentclass[10pt]{article}
\usepackage{amsmath,amssymb,geometry,url}
\geometry{margin=1in}
\setlength{\emergencystretch}{3em}
\title{There is no chain-length threshold for torsion-free quotients\\A disproof of conjecture 00000004287}
\author{gaochengzhecpu (AI-assisted submission)}
\date{3 October 2026}
\begin{document}
\maketitle

\paragraph{The definition and the assertion.}
The problem defines a pure subgroup by the condition that its quotient is torsion-free. We use exactly that definition. The conjecture asserts that unions of chains preserve purity only up to a threshold $\omega$, and that counterexamples exist at every chain length beyond $\omega$. In fact, purity in the stated sense is preserved by every nonempty directed union. In particular, no chain of any ordinal length can be the claimed counterexample.

This argument is valid for abelian groups, the usual setting of the problem, and more generally for arbitrary groups with normal subgroups. Normality is precisely the condition needed to form a quotient group; it is automatic for subgroups of an abelian group.

\paragraph{Quotient criterion.}
Let $G$ be a group and $H$ a normal subgroup. Then
\[
 G/H\text{ is torsion-free}\quad\Longleftrightarrow\quad
 \text{for every }x\in G\text{ and }n\ge1,\quad x^n\in H\Longrightarrow x\in H.
\]
Indeed, $(xH)^n=H$ holds exactly when $x^n\in H$, and $xH=H$ holds exactly when $x\in H$. Every element of the quotient has a representative $x$, so these implications in both directions give the equivalence. In additive abelian notation this criterion is $nx\in H\Rightarrow x\in H$ for every positive integer $n$.

\paragraph{Directed unions.}
Let $(H_i)_{i\in I}$ be a nonempty directed family of normal subgroups of a fixed group $G$, with $G/H_i$ torsion-free for every $i$. Put $U=\bigcup_{i\in I}H_i$.

First, $U$ is a normal subgroup. It contains the identity because the family is nonempty. If $x\in H_i$ and $y\in H_j$, choose $k$ with $H_i,H_j\subseteq H_k$; then $xy\in H_k\subseteq U$. Inverses and conjugates of an element of $H_i$ remain in $H_i$, so these also remain in $U$.

Now suppose $n\ge1$ and $x^n\in U$. There is an index $i$ such that $x^n\in H_i$. The quotient criterion for $H_i$ gives $x\in H_i$, hence $x\in U$. Applying the same criterion to $U$ proves that $G/U$ is torsion-free. This proves the assertion for every nonempty directed family.

Every nonempty chain is directed: among any two of its members, the larger contains both. Therefore the conclusion holds for chains of arbitrary cardinality and arbitrary ordinal length, including every length beyond $\omega$. The claimed counterexamples cannot exist, so the conjecture is false. No assertion about the proposed ranks of those nonexistent counterexamples is needed.

\newpage
\paragraph{What the Lean project proves.}
The accompanying \texttt{lean/Main.lean} is a self-contained Lean 4.19.0 project using only \texttt{Std}. It does not assume commutativity: \texttt{GroupData} records associativity, identity, and both inverse laws, using additive-looking operation names for a general group. A \texttt{Subgroup} records its full membership predicate, subgroup closure, and normality.

The quotient is an actual Lean quotient type, constructed from the coset equivalence relation
\[
 x\sim_H y\quad\Longleftrightarrow\quad \exists h\in H:\ xh=y.
\]
The proof establishes equivalence of this relation, proves compatibility with multiplication and inverse using normality, and constructs all group operations and group laws on the quotient in \path{quotientStructure}. Natural-number multiples in the code mean repeated multiplication in a general group. The lemma \path{project_multiples} proves that the quotient projection preserves them, and \path{project_zero_iff} proves that the kernel of the projection is exactly $H$.

Thus \texttt{Pure} is defined by torsion-freeness of this actual quotient group. It is not defined to be the root-closed criterion. The theorem \path{pure_iff_rootClosed} proves the equivalence between the two notions, and \path{directed_union_pure} then proves closure under directed unions. The union itself is an actual normal subgroup, with its closure and normality proved from the directed family. The theorem \path{arbitrary_chain_union_pure} applies this result to every nonempty chain.

The final theorem \path{conjecture04287_false} rules out the existence of any group with a nonempty chain of pure subgroups whose union is not pure. Any counterexample at a specified transfinite length would be such a chain after forgetting the length. The formal result is therefore stronger than the failure of the proposed threshold; there is no unproved numerical or ordinal cut-off calculation.

The project builds with \texttt{lake build} and \texttt{warningAsError=true}. The final theorem depends only on Lean's standard \texttt{propext} and \texttt{Quot.sound}. There are no proof placeholders, additional axioms, or external decision procedures.
\end{document}
